\documentclass[11pt]{article}

\usepackage{amsmath,amssymb,graphicx}
\usepackage[margin=1in]{geometry}
\usepackage{booktabs}

\title{Unfrozen Kernel Back-Reaction and Mesoscopic Substrate Reinforcement Probes on a Kernel-Induced Cochain Operator Architecture}
\author{Tomislav Rupi\'c \\ Independent Researcher}
\date{March 2026}

\begin{document}

\maketitle

\begin{abstract}
Stage 19 opens Phase II of the HAOS-IIP program by relaxing the strongest Phase I boundary in the smallest possible way: the kernel support and projector remain fixed, but the kernel weights are allowed to become weakly dynamical. The question is narrow and falsifiable. Can bounded substrate back-reaction generate the first measurable historical stabilization effect that the frozen architecture did not support? Five weak substrate laws are tested in sequence on the same representatives inherited from the late pre-geometric atlas: a clustered composite anchor, a phase-ordered symmetric triad, and a counter-propagating corridor. Stage 19A uses local energy-density back-reaction on the kernel edge weights. Stage 19B Branch 1 replaces that local source with mesoscopic basin-residence reinforcement. Stage 19B Branch 2 replaces basin residence with transport-corridor reinforcement built from a smoothed flux-magnitude field with exponential memory. Stage 19C switches symmetry class and introduces bounded flux-driven directional substrate back-reaction. Stage 19D combines weak scalar density memory and weak directional transport memory in a mixed reinforcement law. All five pilots remain bounded, numerically active, and cleanly negative. No composite lifetime gain appears, no binding persistence gain appears, no coarse basin persistence gain appears, no corridor dwell or locking gain appears, no new stabilized ordering class appears, and no refinement promotion is earned. The deformation hierarchy is nevertheless informative: Stage 19A and Stage 19B Branch 2 remain in the $10^{-6}$ to $10^{-5}$ deformation range, Stage 19B Branch 1 and Stage 19D both reach clearly mesoscopic deformation at the $10^{-3}$ to $10^{-2}$ level, and Stage 19C remains bounded but weakly directional. The bounded conclusion is that, within the tested local, mesoscopic, directional, and mixed reinforcement class, substrate memory exists but remains regime-inert. Phase II therefore begins not with emergent capture, but with a sharper architectural limit: weak substrate-law reinforcement is still insufficient to create historical stabilization on this operator architecture.
\end{abstract}

\section{From Phase I to Phase II}

Phase I closed with a clean boundary statement: frozen pre-geometry can organize, but it does not bind. Stages 10 through 18 established that the frozen kernel-induced cochain architecture supports structured morphology, transient composites, weak relational ordering, and bounded memory-like signatures under some extensions, but it does not promote those signatures into robust collective law. Weak effective binding failed. Weak relational feedback failed. Weak phase synchronization failed. Minimal retarded memory failed. The Phase I capstone paper \emph{Limits of Frozen Pre-Geometry} therefore justified a regime shift rather than another local patch.

Stage 19 is that regime shift. The architecture is no longer strictly frozen in its kernel weights, but the support, symmetry, projector, packet family, and operator sector remain fixed. The new hypothesis is deliberately weak: if historical stabilization exists in this framework at all, perhaps the substrate must remember where activity has been. The tested question is not whether geometry emerges. It is whether bounded substrate memory can produce the first regime-level persistence gain.

\section{Stage 19A: Local Energy-Density Back-Reaction}

Stage 19A opens Phase II with the smallest defensible unfrozen-kernel law,
\[
K_{ij}(t)=K^{(0)}_{ij}\bigl(1+\varepsilon S_{ij}(t)\bigr),
\]
where the stress field evolves by a slow relaxation rule driven by local energy density,
\[
S_{ij}(t+\Delta t)=\left(1-\frac{\Delta t}{\tau_{\mathrm{relax}}}\right)S_{ij}(t)+\rho_{ij}(t)-\bar\rho_{ij}(t),
\qquad \rho_{ij}(t)=|A_{ij}(t)|^2.
\]
The update remains support-preserving, symmetric, and bounded by a hard clamp with $\kappa_{\max}=0.1$. No topology change, thresholding, or mixed phase-amplitude driver is introduced.

The pilot uses exactly three representatives and three conditions:
\begin{itemize}
\item null control,
\item very weak back-reaction with $\varepsilon=0.01$,
\item weak back-reaction with $\varepsilon=0.02$.
\end{itemize}
At the base resolution $n=12$, the result is uniform: all nine runs remain in the label ``no kernel back-reaction effect.'' No composite lifetime gain appears, no binding persistence gain appears, no coarse basin persistence gain appears, and no representative earns promotion to the $(12\rightarrow24)$ refinement ladder.

This null is still scientifically useful because the branch is not dead. The kernel deviation norm is nonzero in every non-null run. For the clustered composite anchor it reaches $2.10\times10^{-6}$ at $\varepsilon=0.01$ and $4.20\times10^{-6}$ at $\varepsilon=0.02$, with maximum local deformation reaching $1.21\times10^{-5}$. The phase-ordered symmetric triad and the corridor control show comparable small but nonzero deformation. So Stage 19A is a real negative: local substrate reinforcement enters the dynamics, but it does not move any regime-level observable.

\section{Stage 19B: Mesoscopic Substrate Reinforcement}

Stage 19B keeps the unfrozen-kernel setting but changes the source class. The substrate law becomes mesoscopic rather than microscopic,
\[
K(x,y,t)=K_0(x,y)\bigl(1+\delta S(x,y,t)\bigr),
\]
with bounded low-frequency signals constructed from collective structure. Two branches have been executed so far.

\subsection{Branch 1: Basin-Residence Reinforcement}

Branch 1 reinforces the kernel where the coarse envelope spends long dwell time in a dominant basin. The signal is aligned with the strongest stable structural feature that appeared in the late Phase I atlas: basin sharpening and coarse persistence. The branch is therefore a direct test of whether occupation history can become historical stabilization.

The pilot again uses the same three representatives, the same null/very-weak/weak ladder, and the same base resolution $n=12$. The strengths are $\delta=0.01$ and $0.02$. The result is again cleanly negative. All nine runs remain in ``no mesoscopic substrate effect.'' Composite lifetime, binding persistence, and coarse basin persistence are unchanged. No new stabilized ordering class appears, and no run is promoted.

What changes is the scale of deformation. The clustered composite anchor now reaches a reinforcement-field norm of $0.3902$ at $\delta=0.01$ and $0.3888$ at $\delta=0.02$, with kernel deviation norms of $3.90\times10^{-3}$ and $7.78\times10^{-3}$, and maximum local deformation rising to $4.96\times10^{-2}$. The symmetric triad and corridor control also show clear mesoscopic deformation. So Branch 1 is stronger in kind than 19A, but still regime-inert.

\subsection{Branch 2: Transport-Corridor Reinforcement}

Branch 2 changes the substrate law again. Instead of reinforcing where activity stays, it reinforces where activity repeatedly travels. The corridor field is built from a smoothed local flux-magnitude proxy with single-horizon exponential decay. In the executable pilot the common settings are:
\begin{itemize}
\item corridor estimator: smoothed flux magnitude,
\item corridor memory class: single-horizon exponential decay,
\item corridor memory horizon: $\tau_c=0.25$,
\item smoothing scale: $\sigma=2.0$ grid units,
\item clamp: $0.05$.
\end{itemize}
This is the first true flow-memory test in the atlas.

Again the pilot remains fully negative. All nine runs stay in ``no mesoscopic substrate effect.'' Composite lifetime, binding persistence, and coarse basin persistence remain unchanged. No new stabilized ordering class appears and no corridor-selective transport locking appears. No run is promoted.

The branch is still alive. The clustered composite anchor reaches reinforcement-field norms of approximately $5.83\times10^{-4}$, the symmetric triad reaches $4.44\times10^{-4}$, and the corridor control reaches $5.39\times10^{-4}$. Kernel deviation norms remain in the $10^{-6}$ to $10^{-5}$ range. So the transport-memory field exists and survives numerically, but it is too weak in the tested form to produce a regime-level shift.

\section{Stage 19C: Flux-Driven Directional Substrate Back-Reaction}

Stage 19C changes the symmetry class of the substrate response rather than the source scale. Instead of asking whether the substrate remembers where activity has accumulated, it asks whether the substrate can store a weak directional trace of transport itself. The pilot introduces a bounded antisymmetric flux-driven bias on the same frozen kernel support, with the same three representatives, the same null/very-weak/weak ladder, and the same base resolution $n=12$.

The persistence-first read is fully negative. All nine runs remain in ``no directional substrate effect.'' Composite lifetime, binding persistence, and coarse basin persistence remain flat against null. No representative earns a new stabilized ordering class and no refinement promotion is triggered.

This is still a real probe. The directional field is nonzero in every non-null run. The clustered composite anchor reaches reinforcement-field norms of $1.74\times10^{-4}$ with kernel deviation norms of $3.47\times10^{-6}$ and maximum local deformation of $1.70\times10^{-5}$. The phase-ordered symmetric triad is weaker, with reinforcement-field norms of $6.86\times10^{-5}$ and kernel deviation norms of $1.37\times10^{-6}$. The corridor control reaches reinforcement-field norms of $2.27\times10^{-4}$ with kernel deviation norms of $4.53\times10^{-6}$ and maximum local deformation of $2.51\times10^{-5}$. The corridor case also shows a nonzero curvature proxy, but no corridor dwell gain and no transport locking. So Stage 19C demonstrates bounded directional deformation without persistence consequences.

\section{Stage 19D: Mixed Scalar-Vector Reinforcement}

Stage 19D asks whether the failure of the pure scalar and pure directional branches is caused by symmetry isolation. The mixed branch combines a weak scalar density-memory channel and a weak vector flux-memory channel under the same bounded kernel clamp. The null, balanced-low $(\epsilon_s,\epsilon_v)=(0.01,0.01)$, and balanced-moderate $(0.02,0.02)$ conditions are run on the same three representatives at $n=12$.

The result remains cleanly negative. All nine runs stay in ``no mixed substrate effect.'' No composite lifetime gain appears, no binding persistence gain appears, no coarse basin persistence gain appears, no corridor dwell gain appears, and no transport locking appears. No run earns follow-up refinement.

Again the branch is active rather than dead. The clustered composite anchor reaches mixed-field norms of $6.50\times10^{-3}$ with kernel deviation norms of $6.50\times10^{-3}$ and maximum local deformation of $3.47\times10^{-2}$. The phase-ordered symmetric triad reaches the strongest mixed response, with field norms of $9.42\times10^{-3}$, kernel deviation norms of $9.42\times10^{-3}$, and maximum local deformation of $2.83\times10^{-2}$. The counter-propagating corridor reaches mixed-field norms of $6.70\times10^{-3}$ with maximum local deformation of $3.38\times10^{-2}$. The mixed branch therefore enters the substrate substantially more strongly than the pure directional branch, and at a scale comparable to the stronger basin-residence branch, but still fails to generate any persistence effect.

\section{Comparative Deformation Hierarchy}

The three Stage 19 probes do not differ primarily in their promotion outcome; all are negative. They differ in how strongly the substrate law enters the dynamics before failing. The comparative deformation hierarchy below summarizes the qualitative outcome and deformation scale.

\begin{table}[h]
\centering
\begin{tabular}{@{}llll@{}}
\toprule
Branch & Driver & Representative deformation scale & Regime outcome \\
\midrule
19A & local energy density & $10^{-6}$ to $10^{-5}$ & negative \\
19B-1 & basin residence & $10^{-3}$ to $10^{-2}$ & negative \\
19B-2 & transport corridor & $10^{-6}$ to $10^{-5}$ & negative \\
19C & directional flux bias & $10^{-6}$ to $10^{-5}$ & negative \\
19D & mixed scalar + vector & $10^{-3}$ to $10^{-2}$ & negative \\
\bottomrule
\end{tabular}
\caption{Stage 19 deformation hierarchy. All five branches remain bounded and numerically active, but none produces a persistence gain, corridor locking gain, or new stabilized class.}
\end{table}

This hierarchy matters. It shows that the null is not merely caused by a dead code path or trivial underflow. The substrate can be perturbed substantially at the mesoscopic level, as in Branch 1 and Stage 19D, without generating historical stabilization. At the same time, some laws remain weak in kind: the transport-corridor and directional-flux laws are both active but stay near the $10^{-6}$ to $10^{-5}$ deformation level. The conclusion is therefore stronger than a single-branch failure. Weak reinforcement remains insufficient across multiple symmetry classes.

\section{Interpretation Boundary}

Stage 19 supports three bounded claims.
\begin{enumerate}
\item Minimal unfrozen-kernel laws can be implemented cleanly while preserving support, symmetry, bounded composition, and the late-Phase-I measurement ladder.
\item Those laws do produce nonzero substrate deformation at the tested strengths.
\item Within the tested local, mesoscopic, directional, and mixed reinforcement class, that deformation does not translate into historical stabilization.
\end{enumerate}

Stage 19 does not support stronger claims. It does not support emergent geometry, metric structure, force-law emergence, or universality of a new collective regime. The most it supports is that bounded substrate memory has been introduced and observed, but it remains too weak or too mismatched in class to create persistent capture on this architecture.

\section{Conclusion}

Stage 19 is the opening Phase II result. It relaxes the frozen-kernel boundary in the smallest available ways and lets the architecture answer a clean question: can bounded substrate memory create the first historical stabilization effect? The answer, so far, is no.

Local energy-density back-reaction is active but too weak and too local to matter. Mesoscopic basin-residence reinforcement is substantially stronger but still fails to change any regime-level observable. Mesoscopic transport-corridor reinforcement is active but weaker in effect size and likewise fails to produce persistence gain or transport locking. Flux-driven directional substrate response also remains persistence-inert. Mixed scalar-vector reinforcement enters the substrate more strongly than the pure directional branch and at a scale comparable to the stronger scalar branch, but still produces no lifetime gain, binding gain, basin gain, or corridor gain.

The resulting statement is sharp:

\begin{quote}
Weak substrate-law reinforcement, across local, mesoscopic, directional, and mixed symmetry classes, is insufficient to generate persistence in this architecture at the tested couplings.
\end{quote}

This leaves Phase II in a scientifically strong state. The program has not drifted into parameter forcing. Instead it has shown, repeatedly and under regime-preserving constraints, that historical stabilization is expensive in this framework. The weak reinforcement family is now closed cleanly enough that the next coherent move is not coefficient escalation or another nearby substrate tweak, but a higher-level boundary report or a genuinely new architectural ingredient.

\end{document}
